
    %%%%%%%%%%%%%%%%%%%%%%%%%%%%%%%%%%%%%%%%%%%%%%%%%%%%%%%%%%%%%%%%%%%%%%%%%%%%%%%%
    %%% Classe do documento
    \documentclass[12pt, addpoints]{exam}
    \UseRawInputEncoding

    %%%%%%%%%%%%%%%%%%%%%%%%%%%%%%%%%%%%%%%%%%%%%%%%%%%%%%%%%%%%%%%%%%%%%%%%%%%%%%%%
    % preambulo.tex
    %
    % Arquivo de configuração de packages para uso o LaTeX, conforme minhas
    % preferências e modelos pessoais.
    %
    % Para maiores informações, visite:
    %    https://github.com/abrantesasf/latex
    %
    % NÃO ALTERE SE NÃO SOUBER O QUE ESTÁ FAZENDO!
    %%%%%%%%%%%%%%%%%%%%%%%%%%%%%%%%%%%%%%%%%%%%%%%%%%%%%%%%%%%%%%%%%%%%%%%%%%%%%%%%


    %%%%%%%%%%%%%%%%%%%%%%%%%%%%%%%%%%%%%%%%%%%%%%%%%%%%%%%%%%%%%%%%%%%%%%%%%%%%%%%%
    %%% Estruturas de controle:
    \input{static/utils/controles}

    %%%%%%%%%%%%%%%%%%%%%%%%%%%%%%%%%%%%%%%%%%%%%%%%%%%%%%%%%%%%%%%%%%%%%%%%%%%%%%%%
    %%% Compilação condicional em PDF:
    \input{static/utils/pdf}

    %%%%%%%%%%%%%%%%%%%%%%%%%%%%%%%%%%%%%%%%%%%%%%%%%%%%%%%%%%%%%%%%%%%%%%%%%%%%%%%%
    %%% Configurações de layout da página:
    \input{static/utils/layout}

    %%%%%%%%%%%%%%%%%%%%%%%%%%%%%%%%%%%%%%%%%%%%%%%%%%%%%%%%%%%%%%%%%%%%%%%%%%%%%%%%
    %%% Configurações para autores:
    \input{static/utils/autores}

    %%%%%%%%%%%%%%%%%%%%%%%%%%%%%%%%%%%%%%%%%%%%%%%%%%%%%%%%%%%%%%%%%%%%%%%%%%%%%%%%
    %%% Cabeçalho e rodapé:
    %%% Atenção! É MUITO DIFÍCIL ajustar apropriadamente os running headers
    %%% e running footers da primeira vez. Você pode confiar no LaTeX e deixar que
    %%% ele faça o ajuste automaticamente ou pode quebrar a cabeça e
    %%% tentar melhorar algo que o LaTeX já faz muito bem, mesmo que não fique
    %%% exatamente do jeito que você quer. O que você prefere? Se quiser ajustar
    %%% manualmente, descomente a linha abaixo e faça as configurações no arquivo
    %%% "utils/headings.tex".
    %\input{static/utils/headings}

    %%%%%%%%%%%%%%%%%%%%%%%%%%%%%%%%%%%%%%%%%%%%%%%%%%%%%%%%%%%%%%%%%%%%%%%%%%%%%%%%
    %%% Configuração para as fontes do documento:
    %%% Se você quiser usar fontes próprias ou do sistema, precisará ajustar as
    %%% configurações no arquivo "utils/fontes.tex".
    %%% ATENÇÃO: por padrão eu utilizo XeLaTeX com fontes proprietárias projetadas
    %%% por Matthew Butterick, adquiridas em https://mbtype.com, que não podem ser
    %%% distribuídas devido à licença de utilização. Se você usar XeLaTeX, você
    %%% DEVERÁ OBRIGATORIAMENTE ajustar as fontes no arquivo "utils/fontes.tex".
    \input{static/utils/fontes}

    %%%%%%%%%%%%%%%%%%%%%%%%%%%%%%%%%%%%%%%%%%%%%%%%%%%%%%%%%%%%%%%%%%%%%%%%%%%%%%%%
    %%% Sumário:
    \input{static/utils/toc}

    %%%%%%%%%%%%%%%%%%%%%%%%%%%%%%%%%%%%%%%%%%%%%%%%%%%%%%%%%%%%%%%%%%%%%%%%%%%%%%%%
    %%% Matemática:
    \input{static/utils/matematica}

    %%%%%%%%%%%%%%%%%%%%%%%%%%%%%%%%%%%%%%%%%%%%%%%%%%%%%%%%%%%%%%%%%%%%%%%%%%%%%%%%
    %%% Notas de rodapé:
    \input{static/utils/footnotes}

    %%%%%%%%%%%%%%%%%%%%%%%%%%%%%%%%%%%%%%%%%%%%%%%%%%%%%%%%%%%%%%%%%%%%%%%%%%%%%%%%
    %%% Referências cruzadas, links, citações:
    \input{static/utils/crossrefs}

    %%%%%%%%%%%%%%%%%%%%%%%%%%%%%%%%%%%%%%%%%%%%%%%%%%%%%%%%%%%%%%%%%%%%%%%%%%%%%%%%
    %%% Configurações para os hiperlinks em PDF:
    \input{static/utils/pdfconfig}

    %%%%%%%%%%%%%%%%%%%%%%%%%%%%%%%%%%%%%%%%%%%%%%%%%%%%%%%%%%%%%%%%%%%%%%%%%%%%%%%%
    %%% Glossário e índice remissivo:
    \input{static/utils/glossindex}

    %%%%%%%%%%%%%%%%%%%%%%%%%%%%%%%%%%%%%%%%%%%%%%%%%%%%%%%%%%%%%%%%%%%%%%%%%%%%%%%%
    %%% Referências bibliográficas:
    \input{static/utils/bibliografia}

    %%%%%%%%%%%%%%%%%%%%%%%%%%%%%%%%%%%%%%%%%%%%%%%%%%%%%%%%%%%%%%%%%%%%%%%%%%%%%%%%
    %%% Cores:
    \input{static/utils/cores}

    %%%%%%%%%%%%%%%%%%%%%%%%%%%%%%%%%%%%%%%%%%%%%%%%%%%%%%%%%%%%%%%%%%%%%%%%%%%%%%%%
    %%% Computação:
    \input{static/utils/computacao}

    %%%%%%%%%%%%%%%%%%%%%%%%%%%%%%%%%%%%%%%%%%%%%%%%%%%%%%%%%%%%%%%%%%%%%%%%%%%%%%%%
    %%% Gráficos:
    \input{static/utils/graficos}

    %%%%%%%%%%%%%%%%%%%%%%%%%%%%%%%%%%%%%%%%%%%%%%%%%%%%%%%%%%%%%%%%%%%%%%%%%%%%%%%%
    %%% Ícones:
    \input{static/utils/icones}

    %%%%%%%%%%%%%%%%%%%%%%%%%%%%%%%%%%%%%%%%%%%%%%%%%%%%%%%%%%%%%%%%%%%%%%%%%%%%%%%%
    %%% Blocos:
    \input{static/utils/blocos}

    %%%%%%%%%%%%%%%%%%%%%%%%%%%%%%%%%%%%%%%%%%%%%%%%%%%%%%%%%%%%%%%%%%%%%%%%%%%%%%%%
    %%% Boxes coloridos:
    \input{static/utils/colorbox}

    %%%%%%%%%%%%%%%%%%%%%%%%%%%%%%%%%%%%%%%%%%%%%%%%%%%%%%%%%%%%%%%%%%%%%%%%%%%%%%%%
    %%% Tabelas:
    \input{static/utils/tabelas}

    %%%%%%%%%%%%%%%%%%%%%%%%%%%%%%%%%%%%%%%%%%%%%%%%%%%%%%%%%%%%%%%%%%%%%%%%%%%%%%%%
    %%% Ambientes de listas:
    \input{static/utils/listas}

    %%%%%%%%%%%%%%%%%%%%%%%%%%%%%%%%%%%%%%%%%%%%%%%%%%%%%%%%%%%%%%%%%%%%%%%%%%%%%%%%
    %%% Ambientes floats e similares:
    \input{static/utils/floats}

    %%%%%%%%%%%%%%%%%%%%%%%%%%%%%%%%%%%%%%%%%%%%%%%%%%%%%%%%%%%%%%%%%%%%%%%%%%%%%%%%
    %%% Ajustes para a classe EXAM:
    \input{static/utils/exam}

    %%%%%%%%%%%%%%%%%%%%%%%%%%%%%%%%%%%%%%%%%%%%%%%%%%%%%%%%%%%%%%%%%%%%%%%%%%%%%%%%
    %%% Ajustes para textos:
    \input{static/utils/texto}

    %%%%%%%%%%%%%%%%%%%%%%%%%%%%%%%%%%%%%%%%%%%%%%%%%%%%%%%%%%%%%%%%%%%%%%%%%%%%%%%%
    %%% Ambientes, aliases, comandos e customizações em geral para serem
    %%% utilizadas nos documentos. Defina aqui qualquer customização específica
    %%% para seus documentos!
    \input{static/utils/customizacoes}

    %%%%%%%%%%%%%%%%%%%%%%%%%%%%%%%%%%%%%%%%%%%%%%%%%%%%%%%%%%%%%%%%%%%%%%%%%%%%%%%%
    %%% Ajuste do layout, espaçamento de linhas e etc.:
    \geometry{a4paper, portrait, left=2.0cm, right=2.0cm, top=2.0cm, bottom=2.0cm}
    \singlespacing

    %%%%%%%%%%%%%%%%%%%%%%%%%%%%%%%%%%%%%%%%%%%%%%%%%%%%%%%%%%%%%%%%%%%%%%%%%%%%%%%%
    %%% Configurações de cabeçalho e rodapé (não use fancyhdr pois ocorre conflito):
    \pagestyle{headandfoot}
    \runningheadrule
    \firstpageheader{}{}{}
    \runningheader{Sem disciplina}{}{Avaliação}
    \firstpagefooter{}{}{Página \thepage\ de \numpages}
    \runningfooter{}{}{Página \thepage\ de \numpages}
    \runningfootrule

    %%%%%%%%%%%%%%%%%%%%%%%%%%%%%%%%%%%%%%%%%%%%%%%%%%%%%%%%%%%%%%%%%%%%%%%%%%%%%%%%
    %%% Configurações para as propriedades do PDF:
    \hypersetup{
    hidelinks,           % Comente para web, descomente para impressão
    colorlinks=true,      % True para web, False para impressão
    pdftitle=JoaoOlhaIssoAqui: Avaliação,
    pdfauthor=Matheus Endlich Silveira,
    pdfsubject=Sem disciplina,
    pdfkeywords=['Sem tag'],
    pdfinfo={
        CreationDate={}, % Ex.: D:AAAAMMDDHH24MISS
        ModDate={}       % Ex.: D:AAAAMMDDHH24MISS
    }
    }
    \input{static/utils/pdfconfig}

    %%%%%%%%%%%%%%%%%%%%%%%%%%%%%%%%%%%%%%%%%%%%%%%%%%%%%%%%%%%%%%%%%%%%%%%%%%%%%%%%
    %%% Começa o documento
    \begin{document}

    %%%%%%%%%%%%%%%%%%%%%%%%%%%%%%%%%%%%%%%%%%%%%%%%%%%%%%%%%%%%%%%%%%%%%%%%%%%%%%%%
    %%% Folha de instruções padronizadas da UVV para a prova
    \pdfbookmark[1]{Instruções}{instrucoes}
    \begin{figure}[H]
        \begin{center}
            \includegraphics[scale=0.3]{{static/imgs/logo-uvv.png}}
        \end{center}	
    \end{figure}

    \vspace{-1cm}

    \begin{table}[H]
        \centering
        \begin{tabular}{|llll|}
            \hline
            \multicolumn{2}{|l|}{\textbf{Disciplina:} Sem disciplina} & 
            \multicolumn{1}{l|}{\multirow{3}{*}{\textbf{\makecell{Nota:\\ \,\\ \,}}}} & 
            \multirow{3}{*}{\textbf{\makecell{Coordenador:\\ \,\\ \,}}} \\ 
            \cline{1-2}
            \multicolumn{2}{|l|}{\textbf{Professor:} Matheus Endlich Silveira \hspace{4.72cm} } & 
            \multicolumn{1}{l|}{} & \\ 
            \cline{1-2}
            \multicolumn{2}{|l|}{\textbf{Aluno:}} & 
            \multicolumn{1}{l|}{} & \\ 
            \hline
            \multicolumn{1}{|l|}{\textbf{Turma:} \hspace{6.3cm} } & 
            \multicolumn{1}{l|}{\textbf{Semestre:}} & 
            \multicolumn{2}{l|}{\textbf{Valor:} 10 pontos} \\ 
            \hline
            \multicolumn{1}{|l|}{\textbf{Data:}} & 
            \multicolumn{3}{l|}{\textbf{Avaliação:}} \\ 
            \hline
        \end{tabular}
    \end{table}

    \begin{center}
        \fbox{\setlength\fboxsep{0.3cm} \fbox{\parbox{15.4cm}{
            \begin{center}
                \textbf{INSTRUÇÕES PARA A REALIZAÇÃO DESTA AVALIAÇÃO:}
            \end{center}
            \begin{itemize}[leftmargin=*]
                \item Esta é a primeira avaliação da disciplina Sem disciplina, 
                    e engloba o conteúdo referente Sem tag.
                \item Esta avaliação contém 50 questões objetivas, \textbf{todas obrigatórias}.
                \item \textbf{Leia atentamente cada questão!} As questões objetivas terão uma,
                    e apenas uma única, resposta a ser marcada.
                \item \textbf{Desligue o celular} antes de começar. A avaliação é
                    \textbf{individual} e \textbf{sem consulta}.
                \item As questões podem ser respondidas, na prova, com lápis ou caneta. Ao final
                    da prova \textbf{{VOCÊ DEVERÁ PREENCHER O GABARITO COM CANETA
                    DE COR PRETA}} \textbf{\underline{OU AZUL}} para que seja feita a correção
                    automatizada através da leitura do padrão de respostas marcado no
                    gabarito.
                \item \textbf{CUIDADO ao preencher o gabarito} pois eles são \textbf{INDIVIDUAIS
                    e NOMEADOS} (cada estudante tem um gabarito próprio específico, com um
                    código de identificação único). \textbf{Questões rasuradas são anuladas}
                    pelo software de leitura do gabarito.
                \item Ao final da prova, \textbf{DEVOLVA A PROVA E O GABARITO} para o professor.
                \item Siga todas as normas de \textbf{Integridade Acadêmica} da disciplina.
                    Alunos que forem flagrados com qualquer espécie de ``cola'' ou trocando
                    informações com outros alunos terão suas avaliações recolhidas, as notas
                    zeradas, e a situação será encaminhada para a coordenação para a aplicação
                    das penalidades previstas pela universidade.
                \item Com 90 minutos de avaliação você terá tempo de até 2,min por questão,
                    em média.
                \item Boa avaliação!
            \end{itemize}
            \vspace{2cm}
        }}}
    \end{center}
    \newpage

    %%%%%%%%%%%%%%%%%%%%%%%%%%%%%%%%%%%%%%%%%%%%%%%%%%%%%%%%%%%%%%%%%%%%%%%%%%%%%%%%
    %%% Inicia o bloco de questões:
    \begin{questions}

    \newpage
    \question

Dada a seguinte fórmula (lógica de primeira ordem):\\

$\forall x \ \exists y \ | \ ama(x,y)$ \\

qual das seguintes sentenças em linguagem natural ela representa, considerando que \( ama(x,y) \) representa que \( x \) ama \( y \)?


\begin{choices}
\choice Alguém ama a todos.
\choice Há alguém que todos amam.
\choice Ninguém ama a todos.
\choice Nenhuma das anteriores.
\choice Todos amam alguém.
\end{choices}

\question

Em uma lista circular duplamente encadeada com \( n \) elementos, o espaço ocupado apenas pelos apontadores é (assuma que um apontador ocupa \( p \) bytes):
\begin{choices}
\choice \( 4np \)
\choice \( (np)^2 \)
\choice \( 2np \)
\choice \( np \)
\choice \( 6np \)
\end{choices}

\question

Considere as seguintes tabelas em uma base de dados relacional:\\

Departamento (\underline{CodDepto}, NomeDepto)\\

Empregado (\underline{CodEmp}, NomeEmp, CodDepto, Salario)\\

Considere a seguinte consulta SQL:

\begin{verbatim}

SELECT d.CodDepto, d.NomeDepto, SUM(e.Salario)

FROM Departamento d

JOIN Empregado e ON d.CodDepto = e.CodDepto

GROUP BY d.CodDepto, d.NomeDepto

HAVING COUNT(e.CodEmp) > 2 AND AVG(e.Salario) > 40;

\end{verbatim}



Qual o resultado da consulta?
\begin{choices}
\choice A consulta não retorna nada pois está incorreta.
\choice Para cada departamento que tem mais que dois empregados e cuja média salarial, considerando todos empregados do departamento, exceto os dois primeiros, é maior que 40, obter o código de departamento, seguido do nome do departamento, seguido da soma dos salários dos empregados do departamento.
\choice Para cada departamento que tem mais que dois empregados e cuja média salarial é maior que 40, obter o código de departamento, seguido do nome do departamento, seguido da soma dos salários dos empregados do departamento.
\choice Para cada departamento que tem mais que dois empregados e cuja média salarial é maior que 40, obter um grupo de linhas que contém, para cada empregado do departamento, o código de seu departamento, seguido do nome de seu departamento, seguido da soma dos salários dos empregados do departamento.
\choice Para cada empregado que tem mais que dois departamentos, ambos com média salarial maior que 40, obter o código de departamento, seguido do nome do departamento, seguido da soma dos salários dos empregados do departamento.
\end{choices}

\question

Considere um problema em que são dados 5 objetos com os seguintes pesos e valores:\\

\text{pesos: } $(W_1, W_2, W_3, W_4, W_5) = (6, 10, 9, 5, 12)$\\

\text{valores: } $(P_1, P_2, P_3, P_4, P_5) = (8, 5, 10, 15, 7)$\\

Além disso, é dada uma mochila que suporta até 30 unidades de peso, para transportar os objetos. O objetivo do problema é preencher a mochila de tal forma que o valor total dos objetos a serem transportados seja o maior possível, mas sem exceder o limite de peso suportado pela mochila. Assuma que é permitido colocar fração de um objeto na mochila. Qual das seguintes alternativas corresponde ao valor máximo obtido no preenchimento da mochila:
\begin{choices}
\choice 43.1
\choice 21.5
\choice 30.34
\choice 38.83
\choice 12.2
\end{choices}

\question

Filtro da mediana é:
\begin{choices}
\choice Nenhuma das respostas acima.
\choice Indicado para detectar bordas em imagens.
\choice Indicado para atenuar ruído com preservação de bordas (i.e., rápidas transições de nível em imagens).
\choice Indicado para detectar tonalidades específicas em uma imagem.
\choice Indicado para detectar formas específicas em imagens.
\end{choices}

\question

Professor Mac Sperto propôs o seguinte algoritmo de ordenação, chamado de Super Merge, similar ao \textit{merge sort}: divida o vetor em 4 partes do mesmo tamanho (ao invés de 2, como é feito no \textit{merge sort}). Ordene recursivamente cada uma das partes e depois intercale-as por um procedimento semelhante ao procedimento de intercalação do \textit{merge sort}. Qual das alternativas abaixo é verdadeira?
\begin{choices}
\choice Nenhuma das afirmações acima está correta
\choice Super Merge não está correto. Não é possível ordenar quebrando o vetor em 4 partes
\choice Super Merge está correto, mas consome tempo maior que \( O(\text{merge sort}) \)
\choice Super Merge está correto, mas consome tempo \( O(\text{merge sort}) \)
\choice Super Merge está correto, mas consome tempo menor que \( O(\text{merge sort}) \)
\end{choices}

\question

Uma integração de Sistemas Computacionais formando uma rede, tipicamente é implementada através da instalação de uma Arquitetura de Rede, que é composta de camadas e protocolos, em cada um dos elementos que compõem esta rede. Considere que estações “conversam” quando aplicações de usuários conseguem comunicar-se, sintática e semanticamente, através da Rede de Computadores. Baseados nesta premissa e em todos os conceitos associados à implementação e utilização das redes de computadores podemos afirmar como certo:


\begin{choices}
\choice Computadores com arquiteturas de redes diferentes conseguem “conversar”.
\choice Computadores com arquiteturas de rede parecidas conseguem “conversar”.
\choice Nenhuma delas é uma afirmação correta.
\choice Computadores com arquiteturas diferentes podem “conversar” através de multiplexadores.
\choice Computadores com arquiteturas de redes diferentes podem “conversar” através de um \textit{gateway} ou conversor de protocolos.
\end{choices}

\question

Dentre as definições a seguir, ligadas ao conceito de visões do modelo relacional, qual delas é \textbf{INCORRETA}?


\begin{choices}
\choice Uma visão é útil por representar uma percepção particular do banco de dados, compartilhado por muitos aplicativos.
\choice O gerenciamento de visões envolve a conversão da consulta do usuário sobre as visões para a consulta sobre as relações base.
\choice Uma visão relacional é uma relação virtual, derivada de relações base a partir da especificação de operações da álgebra relacional.
\choice Programas aplicativos do banco de dados podem ser executados sobre visões de relações da base de dados.
\choice Uma visão relacional é uma relação virtual que nunca é materializada.
\end{choices}

\question

Uma prova de vestibular foi elaborada com 25 questões de múltipla escolha com 5 alternativas. O número de candidatos presentes à prova foi 63127. Considere a afirmação: Pelo menos 2 candidatos responderam de modo idêntico as \( k \) primeiras questões da prova. Qual é o maior valor de \( k \) para o qual podemos garantir que a afirmação é verdadeira.


\begin{choices}
\choice 7
\choice 8
\choice 9
\choice 10
\choice 6
\end{choices}

\question

Um banco de dados OLAP (Online Analytical Processing --- Processamento Analítico

Online) é um banco de dados organizado de forma a dar suporte ao Business

Intelligence (BI --- Inteligência de Negócio) que, de forma simplificada, é o

processo de analisar dados para obter informações que podem ser utilizadas na

tomada de decisões comerciais e estratégicas. Em relação aos bancos de dados

OLAP, marque a afirmação correta:
\begin{choices}
\choice São otimizados para o processamento de grande volume de dados históricos e métricas de negócio para a tomada de decisões estratégicas, permitindo e facilitando a análise de dados para a geração de informações e conhecimentos para a tomada de decisões.
\choice São otimizados para o processamento de grande volume de transações diárias, o que permite a análise e a tomada de decisões em tempo real, favorecendo assim a administração empresarial.
\choice São preparados para responder perguntas do tipo ``Quais alunos estão matriculados na turma de Bancos de Dados I neste semestre?'
\choice Geralmente os dados de origem do OLTP são de bancos de dados OLAP.
\choice São otimizados para consulta e relatórios de rotina de dados.
\end{choices}

\question

Suponha que \( T \) seja uma árvore AVL inicialmente vazia, e considere a inserção dos elementos \(10, 20, 30, 5, 15, 2\) em \( T \), nesta ordem. Qual das sequências abaixo corresponde a um percurso de \( T \) em pré-ordem:
\begin{choices}
\choice 30, 20, 15, 10, 5, 2
\choice 15, 10, 5, 2, 20, 30
\choice 10, 5, 2, 20, 15, 30
\choice 20, 10, 5, 2, 15, 30
\choice 2, 5, 10, 15, 20, 30
\end{choices}

\question

Numa prova de múltipla escolha com 10 questões e 4 alternativas, qual a chance (probabilidade) de um aluno apenas "chutando as respostas" conseguir "gabaritar" a prova (acertar todas as questões).


\begin{choices}
\choice \( \frac{1}{2^{20}} \)
\choice \( \frac{1}{4^{20}} \)
\choice \( \frac{1}{10^{4}} \)
\choice \( \frac{1}{4^{15}} \)
\choice \( \frac{1}{10^{8}} \)
\end{choices}

\question

Em relação aos metadados de um banco de dados, assinale a afirmação correta:
\begin{choices}
\choice Como são sempre atualizados de forma automática pelo sistema de gerenciamento de bancos de dados, estão sempre prontos para nos dar uma ``foto' da estrutura do banco de dados em um momento no tempo.
\choice São atualizados pelo SGBD, sempre que ocorre alguma operação de inserção, atualização ou remoção de dados.
\choice Armazenam informações sobre os dados dos usuários, ou seja, são dados de rotina dos usuários, aqueles dados que os usuários precisam para trabalhar no seu dia a dia.
\choice Armazenam informações como os dados de cadastro dos usuários, os valores de compra em uma ordem de compra, as turmas e disciplinas que os professores dão aula, e dados de recursos humanos.
\choice São armazenados no catálogo (ou dicionário de dados) do sistema de gerenciamento de bancos de dados, que fica nas mesmas tabelas de dados dos usuários, permitindo assim que os metadados estejam sempre atualizados.
\end{choices}

\question

Dentre as características do modelo relacional e do modelo de objetos em bancos de dados, qual afirmação é \textbf{INCORRETA}?


\begin{choices}
\choice O modelo de objetos provê uma representação bem próxima de linguagens de programação.
\choice O modelo de objetos possui mais recursos estruturais para a representação de dados que o relacional.
\choice O relacionamento de herança é diretamente representado no modelo relacional.
\choice O modelo de objetos é mais adequado para a representação de tipos abstratos de dados.
\choice O relacionamento binário \( N \times M \) é representado de modo semelhante nos dois modelos.	
\end{choices}

\question

Os protocolos de transporte atribuem a cada serviço um identificador único, o qual é empregado para encaminhar uma requisição de um aplicativo cliente ao processo servidor correto. Nos protocolos de transporte TCP e UDP, como esse identificador se denomina?
\begin{choices}
\choice Endereço IP
\choice Porta
\choice Identificador do processo (PID)
\choice Conexão
\choice Protocolo de aplicação
\end{choices}

\question

Considere as seguintes afirmativas:

\begin{itemize}

    \item O modelo matemático de uma lista é a sequência linear de itens, cuja principal propriedade estrutural é a posição relativa dos elementos dentro da sequência.

    \item A fila e a pilha são consideradas casos especiais da lista.

    \item Numa fila a inserção e a retirada são feitas no mesmo extremo.

    \item Numa lista a inserção e a retirada podem ser feitas em qualquer posição.

    \item Numa pilha apenas a inserção pode ser feita em qualquer posição.

\end{itemize}



Quais são as afirmativas verdadeiras?
\begin{choices}
\choice somente I e III
\choice somente II, III e IV
\choice todas
\choice somente I, II e IV
\choice somente II, IV e V
\end{choices}

\question Assinale a alternativa \textbf{incorreta}:
\begin{choices}
\choice Modelo de dados é a mesma coisa que o diagrama relacional
\choice Um SGBD não tem interface gráfica
\choice No modelo relacional, tudo o que o usuário percebe são tabelas
\choice O MariaDB é um SGBD
\choice O modelo relacional de dados foi criado por Edgar Frank Cood
\end{choices}

\question

Considere o seguinte trecho de programa:\\

1. \( i := 1; \)\\

2. while \( i \leq n \) do\\

   begin\\

3. \( sum := sum + a[i]; \)\\

4. \( i := i + 1; \)\\

   end;\\



Considere que:\\

- \( I \) representa a inicialização da variável \( i := 1 \) na linha 1;\\

- \( T \) representa o teste da linha 2;\\

- \( A \) representa os comandos da linha 3;\\

- \( P \) representa o incremento na linha 4.\\



Qual é a expressão regular que representa todas as sequências de passos possíveis de serem executados por este trecho de programa?
\begin{choices}
\choice \( I(TAP)^* \)
\choice \( I(T(AP)^*)^* \)
\choice \( I(TAP)^+ \)
\choice \( I(T(AP)^*)^+ \)
\choice \( IT^+A^*P^* \)
\end{choices}

\question

Sistemas de processamento de transações, tais como sistemas de reservas aéreas, devem prover um mecanismo que garanta que cada transação não é afetada por outras transações que possam estar ocorrendo ao mesmo tempo. Transações de duas fases obedecem a um protocolo que garante essa atomicidade. Em transações de duas fases:


\begin{choices}
\choice Verifica-se a disponibilidade de todas as travas antes de executar qualquer ação de travamento.
\choice Qualquer objeto corretamente travado deve ser destravado antes que outro objeto possa ser travado.
\choice Todas as operações de leitura ocorrem antes da primeira operação de escrita.
\choice Todas as ações de travamento (\textit{lock}) ocorrem antes da primeira ação de destravamento.
\choice Uma trava compartilhada sobre um objeto deve ser obtida antes de uma trava exclusiva sobre o objeto ser obtida.
\end{choices}

\question

Em um heap com \( n \) vértices existem:
\begin{choices}
\choice exatamente \( \left\lceil \frac{n}{5} \right\rceil \) folhas
\choice não mais que \( \left\lfloor \frac{n}{5} \right\rfloor \) folhas
\choice não menos que \( \frac{2n}{3} \) folhas
\choice exatamente \( \left\lceil \frac{n}{2} \right\rceil \) folhas
\choice aproximadamente \( \log n \) folhas
\end{choices}

\question

Considere as seguintes afirmações sobre redes neurais artificiais: \begin{enumerate}[label=\Roman*.] \item Um perceptron elementar só computa funções linearmente separáveis. \item Não aceitam valores numéricos como entrada. \item O "conhecimento" é representado principalmente através do peso das conexões. \end{enumerate} São corretas:
\begin{choices}
\choice Apenas I e II 
\choice Apenas II e III
\choice Apenas III
\choice I, II e III
\choice Apenas I e III
\end{choices}

\question

O processo de visualização de objetos 3D envolve uma série de passos desde a representação vetorial de um objeto até a exibição da imagem correspondente na tela do computador (pipeline 3D). Selecione a alternativa abaixo que reflete a ordem correta em que esses passos devem ocorrer. 
\begin{choices}
\choice Transformação de câmera, mapeamento para coordenadas de tela, recorte 3D, rasterização, projeção. 
\choice Recorte 3D, transformação de câmera, rasterização, projeção, mapeamento para coordenadas de tela. 
\choice Projeção, transformação de câmera, recorte 3D, mapeamento para coordenadas de tela, rasterização. 
\choice Nenhuma das respostas acima está correta.
\choice Transformação de câmera, recorte 3D, projeção, mapeamento para coordenadas de tela, rasterização. 
\end{choices}

\question

Seu chefe que fazer uma campanha de venda para produtos que custam entre 20 e 30

reais, mas está preocupado porque ouviu de um dos vendedores que diversos

produtos nessa faixa de preço estão com o estoque zerado. Para verificar a

situação seu chefe pediu para que você prepare um relatório que mostre os

produtos que custam entre 20 e 30 reais e que estão com o estoque zerado. Você

preparou o seguinte relatório para seu chefe:

\begin{tcolorbox}

 \begin{verbatim}1  SELECT l.nome AS loja

2       , p.nome AS produto

3       , p.preco_unitario AS preco

4       , e.quantidade AS qtd

5  FROM produtos p

6  INNER JOIN estoques e ON (e.produto_id = p.produto_id)

7  INNER JOIN lojas l ON (e.loja_id = l.loja_id)

8  WHERE (p.preco_unitario >= 20 OR p.preco_unitario <= 30)

9    AND e.quantidade = 0

10 ORDER BY loja, produto;

 \end{verbatim}

\end{tcolorbox}

O seu relatório:
\begin{choices}
\choice Está correto e trará exatamente, os produtos com o estoque zerado e na faixa de preço requisitada pelo seu chefe. Além disso ordena o resultado pelo nome da loja e pelo nome do produto.
\choice Está errado pois há um erro de sintaxe nas linhas 6 e 7.
\choice Está errado pois vai trazer também produtos fora da faixa de preço que o seu chefe quer.
\choice Nenhuma das alternativas anteriores.
\choice Está errado pois a tabela que deveria estar na cláusula FROM é a tabela ``lojas', não a tabela ``produtos'.
\end{choices}

\question

A função abaixo, escrita na linguagem C, quando executada para \( n = 5 \), faz quantas chamadas recursivas (excluindo a primeira chamada da função)? 

\begin{verbatim} int fat (int n) 

{ 

if (n == 1) return n; 

else return (n*fat(n-1)); 

} 

\end{verbatim}
\begin{choices}
\choice 4
\choice 0
\choice 5
\choice 6
\choice 1
\end{choices}

\question

Todos os convidados presentes num jantar tomam chá ou café. Treze convidados bebem café, dez bebem chá e 4 bebem chá e café. Quantas pessoas tem nesse jantar.


\begin{choices}
\choice 10
\choice 19
\choice 15
\choice 27
\choice 23
\end{choices}

\question

Qual dos protocolos abaixo pode ser caracterizado como protocolo de roteamento do tipo estado de enlace?
\begin{choices}
\choice OSPF
\choice ICMP
\choice BGP-4
\choice RIP2
\choice IGMP
\end{choices}

\question

Uma definição lógica, abstrata e auto-contida dos objetos que modelam a

estrutura de armazenamento de dados, dos operadores que modelam o comportamento

e dos demais conceitos que modelam a integridade, e que, juntos, constituem a

máquina abstrata com a qual os usuários interagem é chamada de:


\begin{choices}
\choice Nenhuma das respostas acima
\choice Modelo relacional
\choice Modelo de dados
\choice Banco de dados
\choice Modelo entidade-relacionamento
\end{choices}

\question

Considere a seguinte tabela em uma base de dados relacional (chave primária sublinhada):\\

Tabela1 (CodAluno, CodDisciplina, AnoSemestre, NomeAluno, NomeDisciplina, CodNota, DescricaoNota)\\

Considere as seguintes dependências funcionais:

\begin{itemize}

    \item CodAluno \(\rightarrow\) NomeAluno

    \item CodDisciplina \(\rightarrow\) NomeDisciplina

    \item (CodAluno, CodDisciplina, AnoSemestre) \(\rightarrow\) CodNota

    \item (CodAluno, CodDisciplina, AnoSemestre) \(\rightarrow\) DescricaoNota

    \item CodNota \(\rightarrow\) DescricaoNota

\end{itemize}



Considerando as formas normais, qual das afirmativas abaixo se aplica:
\begin{choices}
\choice A tabela não está na primeira forma normal.
\choice A tabela encontra-se na terceira forma normal, mas não na quarta forma normal.
\choice A tabela encontra-se na segunda forma normal, mas não na terceira forma normal.
\choice A tabela está na quarta forma normal.
\choice A tabela encontra-se na primeira forma normal, mas não na segunda forma normal.
\end{choices}

\question

Se \( \lim_{n \to \infty} \left( \frac{1}{n^2} + \frac{2}{n^2} + \cdots + \frac{n-1}{n^2} \right) = L \) então


\begin{choices}
\choice \( L = 2 \)
\choice \( L = \infty \)
\choice \( L = 1 \)
\choice \( L = 1/2 \)
\choice \( L = 0 \)
\end{choices}

\question

Sobre a hierarquia de Chomsky podemos afirmar que:
\begin{choices}
\choice Uma linguagem que não é regular é livre de contexto
\choice As linguagens reconhecidas por autômatos a pilha são as linguagens regulares
\choice Uma linguagem que é recursivamente enumerável não pode ser uma linguagem regular
\choice As linguagens livres de contexto e as linguagens sensíveis a contexto se excluem
\choice Há linguagens que não são nem livres de contexto nem sensíveis a contexto
\end{choices}

\question

Considere uma cena representada no sistema de referência do universo (SRU), uma \textit{window} definida pelo par de coordenadas (0,0)- (100,100) e uma \textit{viewport} definida pelo par de coordenadas (20,30)- (300,100). Considere ainda que as coordenadas que definem \textit{window} e \textit{viewport} correspondem, respectivamente, aos limites inferior esquerdo e superior direito de ambas. Analise as afirmativas abaixo levando em consideração os conceitos clássicos de \textit{window} e \textit{viewport} e assinale a alternativa correta.



\begin{itemize}

    \item[I)] \textit{Window} e \textit{viewport} estão definidas no SRU.

    \item[II)] No processo de mapeamento desta \textit{window} para esta \textit{viewport} haverá modificação na relação de aspecto.

    \item[III)] O mapeamento da \textit{window} redefinida pelo par de coordenadas (0,0) – (50,50) para a mesma \textit{viewport} (20,30) – (300,100) corresponde a uma operação de \textit{zoom out} sobre o mesmo universo.

\end{itemize}
\begin{choices}
\choice As alternativas I e III são falsas.
\choice As afirmativas II e III são verdadeiras.
\choice As alternativas I e II são falsas.
\choice As alternativas I e II são verdadeiras.
\choice Apenas a afirmativa III é verdadeira.
\end{choices}

\question

Histograma de uma imagem com \( K \) tons de cinza é:
\begin{choices}
\choice Contagem do número de tons de cinza que ocorreram na imagem.
\choice Nenhuma alternativa acima.
\choice Contagem do número de objetos encontrados na imagem.
\choice Contagem do número de vezes que cada um dos \( K \) tons de cinza ocorreu na imagem.
\choice Contagem dos pixels da imagem.
\end{choices}

\question

Um compilador detecta:
\begin{choices}
\choice erros nos resultados gerados pelo programa
\choice erros de sintaxe do programa 
\choice erros aritméticos
\choice erros que podem ocorrer durante a execução do programa
\choice todos os erros citados acima
\end{choices}

\question

Avalie o banco de dados de professores e alunos, mostrado a seguir. Esse banco

de dados pretende armazenar todas as turmas, sendo que um professor pode dar

aula para vários alunos ao mesmo tempo, e um mesmo aluno pode ter aula com

diferentes professores ao mesmo tempo.

\begin{figure}[H]

  \centering

  \caption{Professores, turmas e alunos}

  \label{fig:professores}

  %\fbox{

    \includegraphics[width=0.5\linewidth]{static/imgs/defaul.png}

  %}\\

  %\footnotesize{Fonte: xxxx.}

\end{figure}

Esse banco de dados resolveráo problema de armazenar os dados das turmas?
\begin{choices}
\choice Sim, pois o diagrama mostra claramente que um professor pode ter zero ou mais turmas, e cada aluno pode participar de zero ou mais turmas também.
\choice Não, pois o relacionamento entre 'turmas' e 'professores' está errado: a PK 'cod\_turma' deveria ter ido para a tabela 'professores' como uma FK, mas no diagrama está ao contrário.
\choice Não, pois as cardinalidades e os sentidos dos relacionamentos entre 'professores' e 'turmas', e entre 'turmas' e 'alunos' não cria um relacionamento N:N entre 'professores' e 'alunos'.
\choice Sim, pois as cardinalidades e os sentidos dos relacionamentos entre 'professores' e 'alunos' indica, no final, um relacionamento N:N entre essas duas tabelas.
\choice Não é possível determinar somente avaliando o diagrama exibido, seriam necessárias mais informações do dicionário de dados.
\end{choices}

\question

Dado um vetor \( u \in \mathbb{R}^2 \), \( u = (-3, 4) \), vamos denotar por \( v \) o vetor de \( \mathbb{R}^2 \) que tem tamanho 1 e é ortogonal a \( u \). Então \( v \) pode ser dado por


\begin{choices}
\choice \( \left( -\frac{4}{5}, -\frac{3}{5} \right) \)
\choice \( \left( -\frac{4}{5}, \frac{1}{5} \right) \)
\choice \( \left( -\frac{4}{5}, \frac{2}{5} \right) \)
\choice \( \left( \frac{3}{5}, \frac{4}{5} \right) \)
\choice \( \left( -\frac{4}{5}, \frac{3}{5} \right) \)
\end{choices}

\question

Considere as seguintes afirmações sobre mecanismos de inferência em sistemas baseados em regras.

\begin{enumerate}[label=\Roman*.]

    \item O encadeamento regressivo tem pouca utilidade prática, pois deve partir do possível resultado.

    \item O encadeamento progressivo tanto pode ser em amplitude quanto em profundidade.

    \item Podem trabalhar com informações incertas ou incompletas. 

\end{enumerate}

São corretas:
\begin{choices}
\choice Apenas I e II
\choice Apenas I e III
\choice Apenas II e III
\choice I, II e III
\choice Apenas III
\end{choices}

\question

Assinale a alternativa que contém a ordem correta das cláusulas SQL:


\begin{choices}
\choice SELECT, FROM, WHERE, GROUP BY, JOIN, HAVING, ORDER BY
\choice Nenhuma das alternativas acima
\choice SELECT, FROM, WHERE, JOIN, HAVING, GROUP BY, ORDER BY
\choice SELECT, FROM, JOIN, WHERE, GROUP BY, HAVING, ORDER BY
\choice SELECT, FROM, JOIN, WHERE, GROUP BY, ORDER BY, HAVING
\end{choices}

\question

(Adaptado do ENADE 2005) Avalie o seguite diagrama relacional de um banco de dados:

\begin{figure}[H]

  \centering

  \caption{Banco de dados de peças e fornecedores}

  \label{fig:pecas}

  %\fbox{

    \includegraphics[width=0.5\linewidth]{static/imgs/defaul.png}

  %}\\

  %\footnotesize{Fonte: xxxx.}

\end{figure}

Assinale a opção que apresenta um comando SQL que permite obter uma lista que

contenha o nome de cada fornecedor que tenha fornecido alguma peça, o código da

peça fornecida, a descrição dessa peça e a quantidade fornecida da referida

peça.
\begin{choices}
\choice \begin{verbatim}SELECT * FROM pecas INNER JOIN fornecedores INNER JOIN fornecimentos;\end{verbatim} 
\choice \begin{verbatim}SELECT * FROM pecas INNER JOIN fornecimentos f ON (p.codigo = f.cod_peca) INNER JOIN fornecedores f2 ON (f.cod_forn = f2.cod_forn) ORDER BY f2.nome;\end{verbatim} 
\choice \begin{verbatim}SELECT f2.nome, p.codigo, p.descricao, f.quantidade FROM pecas INNER JOIN fornecimentos f ON (p.codigo = f.cod_peca) INNER JOIN fornecedores f2 ON (f.cod_forn = f2.cod_forn);\end{verbatim} 
\choice \begin{verbatim}SELECT nome, codigo, descricao, quantidade FROM pecas INNER JOIN FORNECEDORES INNER JOIN FORNECIMENTOS;\end{verbatim} 
\choice \begin{verbatim}SELECT DISTINCT p.nome, p.codigo, p.descricao, f.quantidade FROM pecas p INNER JOIN fornecimentos f ON (p.codigo = f.cod_peca);\end{verbatim} 
\end{choices}

\question

Redes semânticas, frames e lógica são formalismos utilizados principalmente em: 
\begin{choices}
\choice redes neurais
\choice IA distribuída
\choice representação de conhecimento
\choice descoberta de conhecimento em bases de dados
\choice inferência em sistemas especialistas
\end{choices}

\question

Você estava ouvindo a discussão de dois colegas, o João e a Maria. Enquanto o

João afirmava que o conceito de banco de dados (BD) é extremamente diferente do

conceito de sistema de gerenciamento de bancos de dados (SGBD), a Maria

discordava e afirmava que o banco de dados nada mais é do que um componente do

SGBD. Nesse momento o professor entrou na sala, percebeu a discussão e falou

que:
\begin{choices}
\choice A Maria estava certa já que o banco de dados é, realmente, um componente interno do sistema de gerenciamento de bancos de dados e não existe de forma independente.
\choice Os dois estavam certos pois o conceito, na prática, depende da situação específica que se está considerando.
\choice Não é possível saber quem estava certo ou errado pois não definiram qual o modelo de dados que estava sendo considerado na discusão.
\choice O João estava certo já que o banco de dados é, realmente, um conceito diferente do conceito do sistema de gerenciamento de bancos de dados, e existe de forma independente.
\choice Os dois estavam errados pois não são conceitos separados nem componentes, são conceitos sinônimos.
\end{choices}

\question

Em relação ao paradigma de programação cliente-servidor. Qual das afirmativas abaixo é FALSA?
\begin{choices}
\choice Um aplicativo cliente é um programa arbitrário que se torna temporariamente um cliente quando for necessário o acesso remoto a um serviço, mas também executa processamento local.
\choice Um aplicativo servidor aceita contato de clientes arbitrários, mas oferece um único serviço. 
\choice Um aplicativo servidor inicia ativamente o contato com clientes 
\choice Um aplicativo servidor é um programa de propósito especial dedicado a fornecer um serviço, mas pode tratar de múltiplos clientes remotos ao mesmo tempo. 
\choice Um aplicativo cliente pode acessar múltiplos serviços quando necessário.
\end{choices}

\question

Um agente SNMP é um aplicativo que é executado:
\begin{choices}
\choice em "firewalls" com o objetivo de proteger acesso a rede
\choice em computadores denominados de gerentes
\choice em um dispositivo de rede
\choice em roteadores com filtragem de pacotes com o objetivo de proteger acesso a rede
\choice a partir de um computador específico para monitorar a rede
\end{choices}

\question

Três atletas \( A \), \( B \) e \( C \) competiram, ao pares, numa corrida de \( d \) metros. Considerando que cada atleta teve o mesmo desempenho (ou seja, a mesma velocidade) ao competir com adversários distintos, e sabendo-se que



\begin{itemize}

    \item \( A \) venceu \( B \) chegando 20 metros à frente

    \item \( B \) venceu \( C \) chegando 10 metros à frente

    \item \( A \) venceu \( C \) chegando 28 metros à frente,

\end{itemize}



podemos afirmar que a corrida tem
\begin{choices}
\choice 200 metros
\choice 110 metros
\choice 100 metros
\choice 50 metros
\choice 150 metros
\end{choices}

\question

Considere as afirmações abaixo sobre endereços IP:

\begin{itemize}

    \item[I.] Uma rede IP classe C fornece até 256 endereços válidos para serem atribuídos a equipe.

    \item[II.] A quantidade máxima de bits que pode ser utilizada para se definir sub-redes em uma rede IP classe C é seis (6).

    \item[III.] A máscara padrão para uma rede classe B é 255.255.255.0.

\end{itemize}

Qual das seguintes alternativas representa as afirmações corretas?
\begin{choices}
\choice Somente II.
\choice Somente I e II.
\choice Somente I.
\choice Somente III.
\choice Somente II e III.
\end{choices}

\question

Qual das afirmações a seguir, relativas à análise sintática, está INCORRETA?


\begin{choices}
\choice Analisadores sintáticos descendentes recursivos são mais simples de implementar do que analisadores sintáticos redutivos.
\choice Uma das diferenças entre os diversos algoritmos de análise redutiva é a forma de identificar o \textit{handle} na pilha.
\choice Algoritmos de análise redutiva podem ser utilizados mesmo para gramáticas ambíguas.
\choice Algoritmos descendentes recursivos podem ser utilizados para algumas gramáticas ambíguas.
\choice As gramáticas LL podem descrever mais linguagens do que as gramáticas LR.
\end{choices}

\question

Considere um grafo \( G \) satisfazendo as seguintes propriedades:



\begin{itemize}

    \item \( G \) é conexo

    \item Se removermos qualquer aresta de \( G \), o grafo obtido é desconexo.

\end{itemize}



Então é correto afirmar que o grafo \( G \) é:
\begin{choices}
\choice Um circuito
\choice Não bipartido
\choice Hamiltoniano
\choice Uma árvore
\choice Euleriano
\end{choices}

\question

Na definição da aquisição de um novo software de

gerenciamento de bancos de dados (SGBD) para uma empresa da área de transporte 

coletivo urbano, a direção da área de informática conduziu o processo de decisão

da seguinte forma: foi designado um profissional da área de bancos de dados

(aquele com maior experiência na área) e atribuída a ele a tarefa de decidir

qual seria o melhor SBGD a ser adquirido. Esse profissional desenvolveu uma

série de estudos sobre as opções disponíveis, utilizando ténicas de simulação e

testes específicos, levando em conta que a quase totalidade dos dados que serão

armazenados são dados de cadastro altamente estruturados e que devem estar

disponíveis na Internet para consulta pelos usuários. Ao final, apresentou

ao diretor um relatório em que indicava claramente qual o melhor tipo de SGBD a

ser adquirido pela empresa. Com base nessa informação, o diretor da empresa

disparou o processo de compra do SGBD. Podemos concluir que o profissional de

bancos de dados escolheu um:
\begin{choices}
\choice Nenhuma das respostas acima.
\choice SGBD Hierárquico, para armazenar corretamente e de forma otimizada todas os relacionamentos entre os ônibus, empregados e passageiros do transporte público.
\choice SGBD Relacional, para armazenar de forma otimizada os dados que, em sua maioria, são altamente estruturados.
\choice SGBD Não Relacional (NoSQL) já que, para armazenar todos os dados e ainda permitir a consulta pelos usuários na Internet, o melhor seria utilizar uma solução com bancos de dados não relacionais que são otimizados para uso na web.
\choice SGBD Grafo, para armazenar todos os relacionamentos entre os usuários do transporte público nos ônibus.
\end{choices}

\question

O contador da figura abaixo é:

\begin{figure}

    \centering

    \includegraphics[width=0.5\linewidth]{static/imgs/defaul.png}

    \label{fig:enter-label}

\end{figure}
\begin{choices}
\choice anisócrono
\choice assíncrono
\choice auto-sincronizado
\choice isócrono
\choice síncrono
\end{choices}

\question

Considerando a rede de Petri abaixo, quais das alternativas são verdadeiras?



\begin{itemize}

    \item[I)] O lugar \( A \) está habilitado a disparar.

    \item[II)] Apenas a transição \( T1 \) está habilitada a disparar.

    \item[III)] A sequência de transições \( (T1, T2, T3, T2) \) pode ser disparada, nessa ordem.

    \item[IV)] A transição \( T4 \) nunca poderá ser disparada.

\end{itemize}



\begin{figure}[H]

    \centering

    \includegraphics[width=0.5\linewidth]{static/imgs/defaul.png}

    \caption{Questão 49}

    \label{fig:enter-label}

\end{figure}
\begin{choices}
\choice Todas as alternativas.
\choice Apenas as alternativas II e III.
\choice Apenas as alternativas II, III e IV.
\choice Apenas as alternativas III, IV.
\choice Apenas as alternativas I e III.
\end{choices}

\question

Um sistema centralizado é um concentrador de recursos; um sistema distribuído apresenta seus recursos dispersos. Entretanto, nem todo conjunto de recursos computacionais dispersos pode ser considerado um sistema distribuído. Considerando um conjunto de computadores, assinale a alternativa que melhor corresponde às características necessárias para considerá-lo um sistema distribuído:
\begin{choices}
\choice Existência de sistema operacional idêntico e hardware padronizado em todos os computadores
\choice Existência de memória compartilhada e relógios locais sincronizados 
\choice Existência de memória secundária compartilhada e protocolos de sincronização de estado
\choice Suporte de rede e funções primitivas de comunicação 
\choice Suporte de rede e um relógio global
\end{choices}



    %%%%%%%%%%%%%%%%%%%%%%%%%%%%%%%%%%%%%%%%%%%%%%%%%%%%%%%%%%%%%%%%%%%%%%%%%%%%%%%%
    %%% Finaliza o bloco de questões:
    \end{questions}
    \end{document}
    